% =====================================================================
%  K-LD :: front matter -- Introduction, Background, Related Work.
%  (Replaces the cloned ESBMC-Arduino body.) The Evaluation is provided by
%  _sec_evaluation.tex (RQ1/RQ4) and _sec_e3.tex (RQ2/RQ3).
% =====================================================================
\providecommand{\kld}{\textsc{K-LD}}
\providecommand{\kst}{\textsc{K-ST}}

\section{Introduction}
\label{sec:intro}

Programmable Logic Controllers run much of the world's safety-critical
automation, and \gls{ld}---the graphical, relay-inspired notation of
\gls{iec}~61131-3---remains the language in which that logic is most often
written and certified. Because a \gls{plc} fault can injure people or damage
plant, \gls{ld} programs are natural targets for formal verification, and a line
of automated tools now checks safety properties of \gls{iec}~61131-3 programs by
translating them into the input language of a model checker. The
\gls{esbmc}-\gls{plc} family, for instance, lowers an \gls{ld} diagram into a
GOTO \gls{ir} and discharges the resulting verification conditions with
\gls{smt}-based \gls{bmc} and $k$-induction.

\paragraph{The trust gap.} Every such tool rests on an unstated assumption: that
its \emph{front-end}---the translation from the diagram to the verifier's model---
faithfully reflects \gls{iec}~61131-3. A verifier is only as sound as this
translation, yet the translation is itself unverified code. In the
\gls{esbmc}-\gls{plc} line the \gls{ld}$\rightarrow$GOTO rules were \emph{designed
to match} an informal reference formalisation, but were never shown equivalent to
the standard. When the translation and the standard disagree, the verifier can
silently return the wrong answer: it can certify a hazardous program safe (if the
translation drops behaviour the standard mandates) or raise false alarms (if it
admits behaviour the standard forbids). Neither failure is detectable from within
the tool, because there is no independent, executable account of what the diagram
\emph{means}.

\paragraph{An executable reference semantics.} We close this gap with \kld{}, an
executable formal semantics of \gls{iec}~61131-3 \gls{ld} defined in the
\textsf{K} framework. \textsf{K} lets one write an operational semantics as a set
of rewrite rules and obtain, for free, both an interpreter (\code{krun}) and a
program verifier (\code{kprove}) from the same definition. \kld{} extends
\kst{}~\cite{Wang2023}, the \textsf{K} semantics of \gls{iec}~61131-3
Structured Text, to the graphical \gls{ld} language: it models normally-open and
normally-closed contacts, energise/latch/unlatch coils, the input--solve--output
\emph{scan cycle} with cross-scan state retention, the on-/off-delay and pulse
timers (\code{TON}, \code{TOF}, \code{TP}), the up/down counters (\code{CTU},
\code{CTD}), and the edge-detection function blocks (\code{R\_TRIG},
\code{F\_TRIG}). Being executable, \kld{} is not a paper artefact: it \emph{runs}
diagrams scan by scan, so its fidelity can be checked directly against a reference
\gls{plc} runtime.

\paragraph{Grounding and using the semantics.} We establish \kld{}'s fidelity and
then put it to work. First, we validate \kld{} scan-for-scan against the reference
function-block implementations that OpenPLC executes---the \gls{matiec}-compiled C
that is the de~facto behavioural standard for these constructs---so that \kld{}'s
timers and counters agree with real \gls{plc} execution on every scan. Second, we
use \kld{} as an \emph{independent reference oracle} to differentially validate
the \gls{esbmc}-\gls{plc} \gls{ld}$\rightarrow$GOTO translation on a suite of
benchmark controllers: for each program we compare the verifier's verdict against
the ground truth obtained by executing the diagram under \kld{}. Third, for the
combinational and latch fragment we go beyond execution and \emph{machine-check},
in \code{kprove}, that \kld{}'s rules implement the \gls{iec}~61131-3
input/output relation for each construct.

The differential study is revealing. On the benchmark suite \kld{} and
\gls{esbmc} agree on the large majority of programs, and on the unsafe variants
\kld{} independently reproduces every violation, localised to the offending
property with a concrete witness. Every \emph{disagreement}, however, is a genuine
\gls{esbmc} defect that \kld{} exposes and OpenPLC confirms---and they fall into
two opposite classes, both rooted in \gls{esbmc}'s incomplete handling of timer
function blocks. On graphical diagrams \gls{esbmc} \emph{skips} the timer path and
thereby \emph{misses a real violation}, certifying as safe a program that
violates its own stated property (an unsound, dangerous error). On the simple
diagram format it instead leaves timer outputs \emph{unconstrained}, manufacturing
counterexamples that no real timer can produce (a false alarm that wastes
engineering effort). An executable, standard-faithful semantics is exactly the
instrument needed to tell these apart from true defects.

\subsection{Contributions}
\label{sub:contributions}
This paper contributes:
\begin{enumerate}
  \item \textbf{\kld{}, an executable formal semantics of \gls{iec}~61131-3
    \gls{ld} in \textsf{K}} (\S\ref{sec:semantics}), extending \kst{} to the
    graphical language: contacts, coils and latches, the retentive scan cycle,
    timers (\code{TON}/\code{TOF}/\code{TP}), counters (\code{CTU}/\code{CTD}),
    and edge function blocks. From one definition we obtain both an interpreter
    and a deductive verifier.
  \item \textbf{Fidelity validation against OpenPLC/\gls{matiec}}: \kld{}'s
    function-block semantics are shown to agree scan-for-scan with the reference
    implementations that OpenPLC executes, grounding \kld{} as a behavioural
    reference for real \gls{plc}s.
  \item \textbf{A differential validation of the \gls{esbmc}-\gls{plc}
    \gls{ld}$\rightarrow$GOTO translation} using \kld{} as an independent oracle
    (\S\ref{sec:e3}), over a benchmark suite covering the full construct fragment
    and both simple and graphical diagram formats. We report broad agreement and
    three corroborated discrepancies, and classify them into \gls{esbmc}'s two
    timer-translation failure modes---unsound \emph{skip} (missed bugs) and
    imprecise \emph{havoc} (false alarms).
  \item \textbf{Machine-checked construct-correctness lemmas} for the
    combinational and latch fragment, discharged in \code{kprove}
    (\S\ref{sec:eval:rq4}), raising the correctness argument for that fragment
    from empirical to formally proven.
\end{enumerate}

\paragraph{Roadmap.} \S\ref{sec:background} reviews \gls{ld} and the scan cycle,
the \textsf{K} framework and \kst{}, the OpenPLC/\gls{matiec} toolchain, and the
\gls{esbmc}-\gls{plc} translation. \S\ref{sec:related} surveys related work.
\S\ref{sec:semantics} presents \kld{} and its validation, \S\ref{sec:e3} the
differential study, and \S\ref{sec:eval:rq4} the mechanised lemmas.

\section{Background}
\label{sec:background}

\subsection{Ladder Diagram and the Scan Cycle}
\label{sub:ld}
An \gls{iec}~61131-3 \gls{ld} program is a ladder of \emph{rungs} drawn between
two power rails. Along a rung, \emph{contacts} test Boolean signals---a
normally-open contact (\code{XIC}) passes power when its variable is true, a
normally-closed contact (\code{XIO}) when it is false---and series and parallel
arrangements of contacts realise conjunction and disjunction. Power reaching the
right rail drives a \emph{coil}: an energise coil (\code{OTE}) copies the rung
condition to its variable, while latch and unlatch coils (\code{OTL}/\code{OTU})
set or reset it and \emph{retain} the value across scans. Rungs may also contain
\emph{function-block} instances---timers (\code{TON}, \code{TOF}, \code{TP}),
counters (\code{CTU}, \code{CTD}), and edge detectors (\code{R\_TRIG},
\code{F\_TRIG})---each carrying private state.

A \gls{plc} executes a program by the cyclic \emph{scan}: it samples the inputs,
solves every rung top to bottom, writes the outputs, and repeats. State that is
not recomputed each scan---latched coils, timer accumulators, counter values---
persists between scans, which is precisely what makes timing and sequencing
behaviour possible and what a faithful semantics must model. Diagrams are
exchanged in the \textbf{PLCopen XML} format, which encodes each rung as a
netlist of contact, coil, and block elements wired by identifier references; the
benchmarks used here are PLCopen XML documents.

\subsection{The \texorpdfstring{\textsf{K}}{K} Framework and \texorpdfstring{\kst{}}{K-ST}}
\label{sub:kframework}
The \textsf{K} framework~\cite{Rosu2010k} is a semantic framework in which the
formal semantics of a language is given as a configuration (a structured
collection of cells) together with rewrite rules over it. From a single
definition \textsf{K} generates an executable interpreter and a reachability-logic
prover (\code{kprove}), so the \emph{same} artefact both runs programs and proves
properties about them---a strong fit for a reference semantics that must be
simultaneously trustworthy and checkable. \textsf{K} has been used to give
complete executable semantics to C, Java, JavaScript, and the EVM, among others.
Closest to our work, Wang et al.~\cite{Wang2023} define \kst{}, a \textsf{K}
semantics for \gls{iec}~61131-3 Structured Text. \kld{} adopts the same framework
and scan-cycle discipline but targets the graphical \gls{ld} language and its
function blocks, which \kst{} does not cover.

\subsection{OpenPLC, \texorpdfstring{\gls{matiec}}{MATIEC}, and Reference Behaviour}
\label{sub:openplc}
OpenPLC is a widely used open-source \gls{plc} runtime; it compiles
\gls{iec}~61131-3 programs to C using \gls{matiec}~\cite{deSousa2014} and
executes them on a scan loop. Crucially, \gls{matiec} ships reference C
implementations of the standard function blocks---the \code{TON}, \code{TOF},
\code{TP}, \code{CTU}, and \code{CTD} bodies---driven by a controlled notion of
elapsed time. Because these implementations are what actually runs on deployed
open-hardware \gls{plc}s, we treat their per-scan behaviour as the behavioural
reference against which \kld{}'s function blocks are validated, and as the
tie-breaker whenever a verifier and \kld{} disagree.

\subsection{Verifying \texorpdfstring{\gls{ld}}{LD}: \texorpdfstring{\gls{esbmc}}{ESBMC}-PLC and the Translation Gap}
\label{sub:esbmcplc}
\gls{esbmc}~\cite{gadelha2020} is a mature \gls{smt}-based \gls{bmc} and
$k$-induction tool for C and other languages. The \gls{esbmc}-\gls{plc} line adds
an \gls{ld} front-end that parses PLCopen XML, lowers each rung into a GOTO
\gls{ir} modelling one scan, wraps it in a scan loop, and checks user-supplied
safety properties. The lowering is engineered to agree with an informal reference
formalisation of \gls{ld}~\cite{Ebnenasir2023}, but there is no proof that it
preserves \gls{iec}~61131-3 semantics. This is the gap \kld{} addresses: it
supplies the missing executable, standard-faithful account of \gls{ld} against
which the lowering can be checked---empirically, by differential testing, and
partially by machine-checked lemmas.

\section{Related Work}
\label{sec:related}

\paragraph{Formal semantics of \gls{iec}~61131-3.} Semantics have been proposed
for several \gls{iec}~61131-3 languages, using timed automata, transition systems,
and abstract state machines, chiefly to enable model checking of Structured Text,
Instruction List, or Sequential Function Charts. \kst{}~\cite{Wang2023} is the
first \textsf{K} semantics for the family and the direct ancestor of this work;
\kld{} extends the framework to the graphical \gls{ld} language and its function
blocks. Unlike semantics built solely to feed a model checker, an executable
\textsf{K} definition doubles as a reference interpreter, which is what makes the
differential validation of another tool possible.

\paragraph{\gls{plc} verification tools.} A range of tools verify
\gls{iec}~61131-3 controllers---among them PLCverif, Arcade.PLC, and SMV/nuXmv-based
flows---each translating the program into a model checker's input. All share the
front-end trust assumption discussed above; none provides an independent
executable semantics to validate that translation. Our contribution is orthogonal
and complementary: rather than proposing another verifier, we provide a reference
oracle with which \emph{any} such translation can be scrutinised, and we apply it
to \gls{esbmc}-\gls{plc}.

\paragraph{Translation validation and semantics in \textsf{K}.} Establishing that
a front-end or compiler preserves semantics is the subject of translation
validation and verified compilation. Full verified compilation of a \gls{plc}
front-end would require a mechanised semantics of both the source diagram and the
verifier's \gls{ir}; we take the lighter, high-coverage route of differential
validation against an executable reference, complemented by machine-checked
per-construct lemmas. Using an executable \textsf{K} semantics both to test and to
prove is enabled by \textsf{K}'s dual interpreter/prover generation, and mirrors
how \textsf{K} semantics of mainstream languages have been used to find and rule
out discrepancies in real implementations.
